%% kapa: exam=1 round=29 year=2018 period=1 form=A subject=economics subject_ko=경제학원론 seq=39 number=79
%% source: Q-Net 감정평가사 기출문제 (cst00309) artlSeq=5202808 file=extracted/1008/5202808/2018년 제29회 감정평가사 1차시험 기출문제/2018년 제29회 감정평가사 1차시험 기출문제/1교시 A형.pdf
%% conversion: kordoc 4.24.51 -> tools/gen1.py (auto); review=proofread-sonnet; flags=none
\begin{question}{79}
총수요-총공급 모형에서 일시적인 음(-)의 총공급 충격이 발생한 경우를 분석한 설명으로 옳지 않은 것은? (단, 총수요곡선은 우하향, 총공급곡선은 우상향한다.)
\begin{choices}
  \item 확장적 통화정책은 국민소득을 감소시킨다.
  \item 스태그플레이션을 발생시킨다.
  \item 단기 총공급곡선을 왼쪽으로 이동시킨다.
  \item 통화정책으로 물가 하락과 국민소득 증가를 동시에 달성할 수 없다.
  \item 재정정책으로 물가 하락과 국민소득 증가를 동시에 달성할 수 없다.
\end{choices}
\end{question}
\answer{1}
